\chapter{Finding II --- the horizon decides the strategy}
\label{ch:horizon}

\section{The puzzle}
The history contains a genuine paradox. The cross-talk network reproduced primate cueing
behaviour on the \emph{long} twenty-nine-step task (hit $\approx0.80$) trained from reward
with no pretrained VAE. The \emph{same} architecture, and the exact paper network, collapse
on the \emph{short} seven-step task without a pretrained VAE --- even though the short task
is, by every naive measure, easier (fewer frames, less accumulated noise, a fixed change
time the agent could memorise). How can the easier task be the one that fails?

\section{Two ways to detect a change}
The resolution is that there are two strategies for change detection, and they have
opposite horizon dependence.
\begin{description}
\item[Instantaneous detection.] Compare \emph{this} frame's orientation to the remembered
baseline. This needs accurate per-frame orientation perception, and it detects the change
in a single frame.
\item[Temporal / flow detection.] Notice that \emph{something shifted} in the visual stream
across frames, without representing what the orientation is. The per-frame signal is weak
and noisy, so it must be integrated over many frames to become reliable; it is slow.
\end{description}

\section{The horizon selects between them}
On the long task the change persists for up to $\sim\!18$ frames, so the flow strategy
works: the agent accumulates a weak ``something moved'' signal until it is confident. It
never \emph{needs} accurate perception, and consequently never learns it --- which is
exactly why the long-trained front-end has orientation decodability $\Rsq\approx0.24$ yet
still solves the long task. On the short task the change lasts two frames. There is no time
to integrate, the flow strategy fails, and the only thing that works is instantaneous
perception --- which sparse reward cannot build (Chapter~\ref{ch:perception}). The
pretrained VAE supplies exactly the missing instantaneous perception, and the short task is
solved.

In one sentence: \textbf{the short task removes the temporal-integration escape hatch and
forces abstract, single-frame perception.} Pixel/flow inference is slow; abstract-but-
accurate inference is fast.

\section{The behavioural fingerprint}
The theory makes a sharp, falsifiable prediction that the data already bear out. A model
that detects instantaneously should press the moment the change appears, so its reaction
time should be pinned at change onset and \emph{cannot} be sped up by the cue. The
pretrained-VAE model on the short task does exactly this: median reaction time is five
frames (change onset), essentially flat across change magnitude and across cue validity
(Chapter~\ref{ch:deepdive}). A flow-based model, by contrast, should have long, magnitude-
dependent reaction times. The prediction to complete the argument is a direct comparison of
reaction-time distributions: the from-scratch long-task model should press \emph{late} and
the pretrained-VAE model \emph{early}, the clean signature of ``flow is slow, perception is
fast.''

\section{Consequences for the program}
Two consequences follow. First, \textbf{the short task is the more discriminating test
bed}: because it forecloses the flow shortcut, it forces the network to actually perceive,
and so it exposes architectural differences that the long task hides behind temporal
integration. This is why the cue-attention investigation of Chapter~\ref{ch:attention} is
run on the short task. Second, \textbf{horizon is a design variable, not a nuisance}:
choosing the seven-step task with the frozen VAE gives the cleanest, most mechanism-
revealing setting, while the twenty-nine-step task is the right venue for studying the
\emph{maintenance} of attention across a delay. The empirical paper will want both: the
short task for the sensitivity/criterion story and the long task for the maintenance and
anticipation story.
